\chapter{GIỚI THIỆU ĐỀ TÀI}

\section{Lý do chọn đề tài}

Trong bối cảnh chuyển đổi số và sự phát triển mạnh mẽ của công nghệ thông tin, phần mềm ngày càng
đóng vai trò quan trọng trong hầu hết các lĩnh vực của đời sống, từ tài chính, y tế, thương mại
điện tử đến các hệ thống hạ tầng và dịch vụ công nghệ. Cùng với sự gia tăng về quy mô và độ phức
tạp của phần mềm, các vấn đề liên quan đến an toàn và bảo mật cũng ngày càng trở nên đáng quan tâm.
Trong đó, lỗ hổng phần mềm (Software Vulnerability) là một trong những nguyên nhân chính dẫn đến
các cuộc tấn công mạng, rò rỉ dữ liệu và gây ảnh hưởng nghiêm trọng đến tính bảo mật, toàn vẹn và
khả dụng của hệ thống. Do đó, việc nghiên cứu và phát triển các phương pháp \ac{PHLHPM} có ý nghĩa
quan trọng trong việc nâng cao chất lượng và đảm bảo an toàn cho các sản phẩm phần mềm.

Các phương pháp \ac{PHLHPM} truyền thống thường dựa trên phân tích tĩnh (Static Analysis), phân
tích động (Dynamic Analysis) hoặc các tập luật và mẫu nhận diện được định nghĩa trước. Mặc dù những
phương pháp này đã đạt được nhiều kết quả nhất định, chúng vẫn tồn tại một số hạn chế trong việc
phát hiện các lỗ hổng phức tạp, đặc biệt khi mã nguồn có cấu trúc đa dạng và liên tục thay đổi.
Trong những năm gần đây, sự phát triển của học máy \ac{ML} và học sâu \ac{DL} đã mở ra hướng tiếp
cận mới cho bài toán \ac{PHLHPM}. Thay vì phụ thuộc hoàn toàn vào các luật được xây dựng thủ công,
các mô hình học máy có thể học những đặc trưng từ mã nguồn để nhận diện các mẫu có khả năng liên
quan đến lỗ hổng bảo mật. Tuy nhiên, việc xây dựng các mô hình này thường đòi hỏi tập dữ liệu lớn
và đa dạng, trong khi dữ liệu mã nguồn thực tế thường được lưu trữ phân tán tại nhiều tổ chức và có
thể chứa các thông tin độc quyền hoặc nhạy cảm. Việc tập trung dữ liệu từ nhiều nguồn về một máy
chủ duy nhất để huấn luyện mô hình có thể phát sinh những rủi ro về quyền riêng tư, bảo mật và
quyền sở hữu trí tuệ, đồng thời gặp phải các hạn chế về chính sách chia sẻ dữ liệu.

\ac{FL}, hay học liên kết, là một phương pháp học máy phân tán cho phép nhiều bên cùng tham gia huấn luyện một mô hình chung mà không cần chia sẻ trực tiếp dữ liệu gốc. Thay vì
truyền dữ liệu đến một máy chủ trung tâm, mỗi bên (client) thực hiện quá trình huấn luyện trên dữ
liệu cục bộ và chỉ gửi các bản cập nhật mô hình (model updates) đến máy chủ điều phối (server) để
tổng hợp. Nhờ đó, \ac{FL} tạo điều kiện cho các tổ chức hợp tác xây dựng mô hình \ac{PHLHPM} mà vẫn duy trì dữ liệu trong phạm vi quản lý của mình. Đây là một hướng tiếp
cận phù hợp đối với những bài toán mà dữ liệu được phân tán và khó chia sẻ trực tiếp giữa các bên.

Mặc dù có ưu điểm trong việc hạn chế trao đổi dữ liệu gốc, \ac{FL} vẫn tồn tại những thách thức
liên quan đến tính minh bạch, khả năng giám sát và truy vết quá trình huấn luyện. Trong một hệ
thống FL, nhiều client có thể tham gia vào các vòng huấn luyện (training rounds) khác nhau, thực
hiện huấn luyện trên dữ liệu riêng và gửi các bản cập nhật mô hình đến server. Tuy nhiên, việc ghi
nhận và xác minh lịch sử đóng góp của từng client không phải lúc nào cũng được hỗ trợ đầy đủ. Khi
mô hình toàn cục (global model) có những thay đổi bất thường về hiệu suất hoặc xuất hiện các kết
quả dự đoán không mong muốn, việc xác định client nào đã tham gia, bản cập nhật nào đã được tiếp
nhận và những cập nhật đó đã ảnh hưởng như thế nào đến mô hình toàn cục có thể gặp nhiều khó khăn.

Bên cạnh đó, môi trường \ac{FL} không thể mặc định rằng tất cả các thành phần tham gia đều đáng tin
cậy. Các client có thể gửi những bản cập nhật mô hình không hợp lệ hoặc bị can thiệp, trong khi máy
chủ điều phối cũng có thể trở thành đối tượng cần được giám sát. Nếu thông tin về quá trình huấn
luyện chỉ được lưu trữ trong cơ sở dữ liệu hoặc hệ thống nhật ký do một thành phần trung tâm quản
lý, các bản ghi này vẫn có nguy cơ bị sửa đổi hoặc xóa bỏ bởi những thực thể có quyền truy cập phù
hợp. Điều này làm hạn chế khả năng kiểm chứng lịch sử hoạt động và xác định nguồn gốc của các bản
cập nhật khi xảy ra sự cố. Vì vậy, bên cạnh việc đảm bảo khả năng huấn luyện phân tán, việc xây
dựng một cơ chế ghi nhận và truy vết đáng tin cậy cho quá trình \ac{FL} cũng là một vấn đề cần được
quan tâm.

Blockchain là một công nghệ tiềm năng có thể hỗ trợ giải quyết vấn đề trên thông qua cơ chế lưu trữ
dữ liệu phân tán, hàm băm mật mã và cấu trúc liên kết các khối dữ liệu. Các thông tin được ghi nhận
trên Blockchain có khả năng kiểm chứng tính toàn vẹn và khó bị thay đổi trái phép khi hệ thống được
triển khai với các cơ chế bảo mật và đồng thuận phù hợp. Khi kết hợp với \ac{FL}, Blockchain có thể
đóng vai trò như một lớp ghi nhận và kiểm toán các sự kiện quan trọng trong quá trình huấn luyện,
chẳng hạn như định danh client, thời điểm tham gia, vòng huấn luyện, mã băm của bản cập nhật mô
hình và trạng thái xử lý. Thông qua đó, hệ thống có thể hỗ trợ truy xuất lịch sử đóng góp của từng
client, xác minh tính toàn vẹn của các thông tin đã ghi nhận và tăng cường khả năng kiểm toán quá
trình huấn luyện. Tuy nhiên, Blockchain không trực tiếp đảm bảo rằng các bản cập nhật được ghi nhận
là hợp lệ hoặc không chứa hành vi độc hại; do đó, việc kết hợp cơ chế ghi nhận với các phương pháp
kiểm tra và xác minh phù hợp là cần thiết.

Mặc dù đã có nhiều nghiên cứu khai thác \ac{FL} cho các bài toán học máy phân tán và ứng dụng
Blockchain nhằm tăng cường tính bảo mật, minh bạch của hệ thống, việc kết hợp các công nghệ này
trong một hệ thống \ac{PHLHPM} vẫn đặt ra những vấn đề cần được nghiên cứu và đánh giá. Đặc biệt,
cần xem xét cách thức liên kết thông tin định danh của client với các sự kiện huấn luyện, cách ghi
nhận và xác minh các bản cập nhật mô hình, cũng như mức độ ảnh hưởng của việc tích hợp Blockchain
đến hiệu năng và chi phí vận hành của hệ thống. Đây là những khía cạnh quan trọng nhằm đánh giá
tính khả thi của việc sử dụng Blockchain để tăng cường khả năng truy vết trong môi trường \ac{FL}.

Xuất phát từ những vấn đề trên, đề tài “Nghiên cứu và triển khai hệ thống \ac{PHLHPM} sử dụng học
liên kết và chuỗi khối” được thực hiện nhằm nghiên cứu và xây dựng một hệ thống \ac{PHLHPM} dựa
trên phương pháp \ac{FL}, kết hợp công nghệ Blockchain để tăng cường tính minh bạch, khả năng truy
vết và kiểm chứng lịch sử huấn luyện. Hệ thống hướng đến việc cho phép nhiều client cùng tham gia
huấn luyện mô hình trên dữ liệu cục bộ, đồng thời ghi nhận các sự kiện quan trọng trong quá trình
huấn luyện nhằm hỗ trợ xác định nguồn gốc và kiểm tra tính toàn vẹn của các bản cập nhật mô hình.
Qua đó, đề tài không chỉ nghiên cứu khả năng ứng dụng học liên kết vào bài toán \ac{PHLHPM} mà còn
đánh giá tiềm năng của Blockchain trong việc nâng cao độ tin cậy và khả năng kiểm toán của môi
trường huấn luyện học máy phân tán.

\section{Mục tiêu nghiên cứu}

\subsection{Mục tiêu tổng quát}

Mục tiêu tổng quát của đề tài là nghiên cứu và triển khai hệ thống \ac{PHLHPM} sử dụng học liên kết
(\ac{FL}) kết hợp với công nghệ chuỗi khối (Blockchain). Hệ thống hướng đến việc xây dựng mô hình
học máy có khả năng phát hiện lỗ hổng trong mã nguồn phần mềm thông qua quá trình huấn luyện phân
tán tại nhiều máy khách, đồng thời ứng dụng Blockchain để ghi nhận và hỗ trợ kiểm chứng các hoạt
động quan trọng trong quá trình huấn luyện.

Thông qua sự kết hợp này, đề tài hướng đến việc hạn chế nhu cầu tập trung mã nguồn từ nhiều bên về
một máy chủ duy nhất, đồng thời tăng cường khả năng truy xuất nguồn gốc, tính toàn vẹn và tính minh
bạch của lịch sử huấn luyện. Từ đó, hệ thống có thể cung cấp cơ sở để theo dõi sự tham gia của các
máy khách và kiểm tra tính nhất quán của các thông tin liên quan đến quá trình cập nhật mô hình.

\subsection{Mục tiêu cụ thể}

Để đạt được mục tiêu tổng quát, đề tài tập trung vào các mục tiêu cụ thể sau:

\begin{itemize}

	\item \textbf{ Nghiên cứu cơ sở lý thuyết và các công trình liên quan: } Tìm hiểu bài toán \ac{PHLHPM}, các phương pháp ứng dụng học máy và học sâu trong phân tích mã nguồn, cơ chế huấn luyện của \ac{FL}, cũng như nguyên lý và khả năng ứng dụng Blockchain trong việc ghi nhận và kiểm chứng dữ liệu.

	\item \textbf{ Xây dựng mô hình \ac{PHLHPM}: } Nghiên cứu và lựa chọn bộ dữ liệu, phương pháp biểu diễn mã nguồn và mô hình học máy hoặc học sâu phù hợp với bài toán \ac{PHLHPM}. Đánh giá khả năng phát hiện của mô hình thông qua các chỉ số thực nghiệm phù hợp.

	\item \textbf{ Triển khai cơ chế học liên kết: } Xây dựng môi trường huấn luyện phân tán gồm máy chủ điều phối và nhiều máy khách. Mỗi máy khách thực hiện huấn luyện mô hình trên dữ liệu cục bộ, sau đó gửi các tham số hoặc cập nhật mô hình đến máy chủ để tổng hợp thành mô hình toàn cục mà không cần chia sẻ trực tiếp mã nguồn huấn luyện.

	\item \textbf{ Tích hợp Blockchain vào quá trình huấn luyện: } Thiết kế cơ chế ghi nhận các sự kiện quan trọng, bao gồm thông tin định danh máy khách, thời điểm tham gia, vòng huấn luyện, trạng thái xử lý và giá trị băm đại diện cho bản cập nhật mô hình. Cơ chế này nhằm hỗ trợ truy xuất lịch sử huấn luyện và kiểm tra tính toàn vẹn của các bản cập nhật đã được ghi nhận.

	\item \textbf{ Xây dựng hệ thống tích hợp và giao diện tương tác: } Kết hợp mô hình \ac{PHLHPM}, quy trình học liên kết và cơ chế ghi nhận trên Blockchain thành một hệ thống thống nhất. Xây dựng giao diện cho phép người dùng theo dõi quá trình huấn luyện, trạng thái tham gia của các máy khách và tra cứu thông tin lịch sử được ghi nhận.

	\item \textbf{ Thực nghiệm và đánh giá hệ thống: } Tiến hành thực nghiệm để đánh giá hiệu quả \ac{PHLHPM} của mô hình, khả năng vận hành của quy trình học liên kết và hiệu quả của cơ chế ghi nhận trên Blockchain. Đồng thời, phân tích chi phí bổ sung do việc tích hợp Blockchain gây ra, chẳng hạn như độ trễ ghi nhận, thời gian xác minh và tài nguyên sử dụng.

\end{itemize}
